\subsection{代数的张量积的根}

设 \(A_{1}\) 与 \(A_{2}\) 是 K-代数，令 \(A = A_{1} \otimes A_{2}\)。

命题 3. —— 假设代数 \(\mathrm{A}_1\) 与 \(\mathrm{A}_2\) 是半单的，其中心分别为 \(\mathbf{Z}_1\) 与 \(\mathbf{Z}_2\)。令 \(\mathbf{Z} = \mathbf{Z}_1 \otimes \mathbf{Z}_2\)。

\begin{enumerate}
\def\labelenumi{\alph{enumi})}
\item
  映射 \(a \mapsto aA\) 是从按包含关系排序的 Z 的理想所成之集到按包含关系排序的 A 的双边理想所成之集的同构。
\item
  A 的根等于 A 的诸极大双边理想之交，且等于 \(\Re(Z)A\)。
\item
  若 K-代数 \(Z_{1}\) 与 \(Z_{2}\) 之一是可分的，特别地若体 K 是完备的，则环 Z 与 A 的根均约化为 0。
\end{enumerate}

每个代数 \(A_{i}\) 都是有限多个单代数的积。而环的积的中心是诸中心的积，且对根（VIII, 第 156 页，推论 3）与对双边理想（I, §8, No.~10, 第 109 页，命题 8）也有类似的断言。因此只需在 \(A_{1}\) 与 \(A_{2}\) 都是单代数的假设下证明命题 3。

对 \(i \in \{1,2\}\)，令 \(B_i = A_i \otimes A_i^o\)，并把 \(A_i\) 视为一个 \(B_i\)-模，其中位似由公式

\[
(x \otimes y) z = x z y
\]

给出，其中 x、y 与 z 属于 \(A_{i}\)。\(B_{i}\)-模 \(A_{i}\) 的交换子是 \((\mathrm{Z}_{i})_{\mathrm{A}_{i}}\)，即由 \(Z_{i}\) 的元素作的位似所成之集；我们把它与 \(Z_{i}\) 视为同一。此外，\(B_{i}\)-子模（\(A_{i}\) 的）就是 \(A_{i}\) 的双边理想，且由于环 \(A_{i}\) 是单环，仅有的双边理想是 0 与 \(A_{i}\)。因而 \(A_{i}\) 是单 \(B_{i}\)-模。再者，\((\mathrm{B}_{1}\otimes\mathrm{B}_{2})\)-子模（\(A_{1}\otimes A_{2}\) 的）就是环 \(A_{1}\otimes A_{2}\) 的双边理想。

于是，把 VIII, 第 214 页，定理 2, b) 应用于单 \(\mathrm{B}_1\)-模 \(\mathrm{A}_1\)（交换子为 \(\mathrm{Z}_1\)）以及单 \(\mathrm{B}_2\)-模 \(\mathrm{A}_2\)（交换子为 \(\mathrm{Z}_2\)），即得断言 a)。

我们来证明断言 b)。A 的诸极大双边理想之交是 \((\mathrm{B}_{1}\otimes\mathrm{B}_{2})\)-模 \(A_{1}\otimes A_{2}\) 的根；由 VIII, 第 215 页，推论 2，这一交与 \(\Re(\mathrm{Z})\mathrm{A}\) 重合。代数 \(Z_{1}\) 与 \(Z_{2}\) 是交换且半单的。因此环 Z 的根由幂零元素组成（VIII, 第 215 页，定理 3），且环 A 的双边理想 \(\Re(\mathrm{Z})\mathrm{A}\) 含于根 \(\Re(\mathrm{A})\)（VIII, 第 157 页，注 1）。然而，A 的诸极大双边理想之交包含 \(\Re(\mathrm{A})\)（VIII, 第 155 页，命题 5, d)）。这就证明了断言 b)。

可分交换代数与既约交换代数的张量积是既约环（V, §15, No.~2, 第 120 页，命题 5）。代数 \(Z_1\) 与 \(Z_2\) 是交换且半单的，因而是既约的。若代数 \(Z_1\) 与 \(Z_2\) 之一是可分的，则代数 \(Z\) 是既约的；于是由 VIII, 第 215 页，定理 3，它是无根的，且我们有 \(\Re(A) = \Re(Z)A = 0\)。特别地，当体 \(K\) 完备时即属此种情形，因为完备体上的每个既约交换代数都是可分的（V, §15, No.~5, 第 125 页，定理 3）。

推论. —— 假设代数 \(A_{1}\) 与 \(A_{2}\) 都是单的，且中心 \(Z_{1}\)（\(A_{1}\) 的）约化为 K。则环 \(A_{1} \otimes A_{2}\) 除 0 与其自身外没有其他双边理想。

由假设，\(Z_{1} = K\)，且由于 \(A_{2}\) 是单的，其中心 \(Z_{2}\) 是一个体（VIII, 第 121 页，推论 1, a)）。因此环 \(Z = Z_{1} \otimes Z_{2}\) 是体，推论由命题 3, a) 得出。

